\documentclass[12pt]{article}
\usepackage[T2A]{fontenc}
\usepackage[utf8]{inputenc}
\usepackage[russian]{babel}
\usepackage{latexsym,amssymb,amsthm}
\usepackage{amsmath}
\usepackage{wrapfig}
\RequirePackage[pdftex]{graphicx}
%%\usepackage{array}

\usepackage{epsf,epic,eepic}
\usepackage{xcolor}
\input ris.tex
\input cells.sty
% \addrisoption{\risleftfalse}

% \font\Chess=chess10 scaled 600

\DeclareRobustCommand{\No}{\ifmmode{\nfss@text{\textnumero}}\else\textnumero\fi}

\oddsidemargin=0pt
\textwidth=159mm
\headheight=0pt
\headsep=0pt
\topmargin=0pt
\textheight=240mm

%%
%% Три точки для обозначения делимости
%%
\catcode`\@=11
\def\kratno{\mathrel{\smash{\lower.5ex\hbox{$\,\vdots\,$}}}}
% Это прямые три точки для обозначения делимости
\def\nekratno{\mathrel{\mathpalette\c@ncel\kratno}}
\def\c@ncel#1#2{\m@th\ooalign{$\hfil#1\mkern1mu/\hfil$\crcr$#1#2$}}
% Это перечеркнутые три точки для обозначения неделимости
\def\divis{\smash{\lower.1ex\hbox{\,\hbox{.}\kern-.22em\lower-.6ex\hbox{.}
\kern-.52em\lower-1.2ex\hbox{.}\,}} }
% А это наклонные...

\hyphenation{чис-ло чис-ла чис-лу чис-лом чис-ле чис-лам чис-ла-ми чис-лах
че-ты-рех-уголь-ник че-ты-рех-уголь-ни-ка пря-мо-уголь-ник тре-уголь-ни-ка
тре-уголь-ни-ков вне-впи-сан-ной вне-впи-сан-ных}

\sloppy
\pagestyle{empty}

\def\q#1.{{\bf #1. }}
\def\dotline{\smallskip\hbox to \hsize{\dotfill}\medskip}

\begin{document}

\centerline{\sc САНКТ-ПЕТЕРБУРГСКАЯ}
\centerline{\sc ОЛИМПИАДА ШКОЛЬНИКОВ 2024 года ПО МАТЕМАТИКЕ}
\centerline{\sc II тур. \ 6 класс.}
\medskip
\hrule
\bigskip

\def\blsq{{\color{lightgray}$\blacksquare$}}

\newbox\kk

\setbox\kk\vbox{\cellsize=8pt
\cput(1,5){\blsq}\cput(2,5){\blsq}\cput(3,5){\blsq}\cput(4,5){\blsq}\cput(5,5){\blsq}
\cput(6,5){\blsq}\cput(7,5){\blsq}\cput(8,5){\blsq}\cput(9,5){\blsq}
\cput(2,2){\blsq}\cput(2,8){\blsq}\cput(8,2){\blsq}\cput(8,8){\blsq}
\cput(3,3){\blsq}\cput(3,7){\blsq}\cput(7,3){\blsq}\cput(7,7){\blsq}
\cput(4,4){\blsq}\cput(4,6){\blsq}\cput(6,4){\blsq}\cput(6,6){\blsq}
\cput(5,1){\blsq}\cput(5,2){\blsq}\cput(5,3){\blsq}\cput(5,4){\blsq}
\cput(5,6){\blsq}\cput(5,7){\blsq}\cput(5,8){\blsq}\cput(5,9){\blsq}\cput(5.05,5){$\bullet$}
\cells{
         _
   _    |_|    _
  |_|_  |_|  _|_|
    |_|_|_|_|_|
 _ _ _|_|_|_|_ _ _
|_|_|_|_|_|_|_|_|_|
     _|_|_|_|_
   _|_| |_| |_|_
  |_|   |_|   |_|
        |_|}}

\q1. В Непале все горы имеют разные названия, разные высоты и все расстояния между ними различны. Два альпиниста пошли покорять вершины. Каждый отправился к какой-то вершине, а затем действовал по следующему алгоритму: покорив вершину, он выбирал из всех более высоких вершин самую близкую и покорял её. Если же более высоких вершин больше не оставалось, альпинист уезжал домой. В результате первый альпинист покорил все вершины, кроме горы Ангелов, а второй~--- все, кроме горы Бесов. Что выше~--- гора Ангелов или гора Драконов?

\q2. Турист приехал на остров, где живут 98 человек, каждый из которых либо рыцарь (всегда говорящий правду), либо лжец (который всегда лжёт). Турист может выбрать любую компанию из пяти человек и спросить у одного из них: \emph{Правда ли, что среди остальных четверых количество рыцарей четно?} Сможет ли турист за 20 вопросов определить, чётное или нечётное количество рыцарей живет на острове?

\setrisbox\copy\kk

\ris
\q3. Снежинка --- это фигура из 29 клеток, изображенная на рисунке (отмечена центральная клетка снежинки). На клетчатой плоскости конечное число клеток покрашено в красный, синий и зеленый цвета. Может ли оказаться так, что в каждой снежинке с~красным центром синих клеток больше, 
чем зеленых, в каждой снежинке с синим центром зеленых больше, чем красных, а в каждой снежинке с зеленым центром красных больше, чем синих? 
\endris

\q4. На доске написано 18 различных натуральных чисел. Докажите, что не могло так получиться, что любое число от 1 до 147 оказалось наибольшим общим делителем каких-то двух (разных) чисел на доске.

\dotline

\centerline{Олимпиада 2024 года. II тур. 6 класс. Выводная аудитория.}
\smallskip

\q5. На кольцевом шоссе расположены 202 деревни. По шоссе ходят 
автобусы, каждый автобус ездит от какой-то одной деревни до какой-то 
другой и обратно тем же путем, останавливаясь также и~во всех деревнях по дороге. 
При этом в каждой деревне останавливается хотя бы один автобус, и для любых 
двух деревень существует автобус, который останавливается ровно в одной из 
них. Какое наименьшее количество автобусов может ходить по шоссе?

\q6. Петя возвел число 4 в некоторую натуральную степень и получил число,
состоящее из четного количества цифр. Вася записал цифры этого числа в 
обратном порядке и обнаружил, что получилась степень числа~5. Докажите, 
что кто-то из мальчиков ошибся.


\clearpage

\centerline{\sc САНКТ-ПЕТЕРБУРГСКАЯ}
\centerline{\sc ОЛИМПИАДА ШКОЛЬНИКОВ 2024 года ПО МАТЕМАТИКЕ}
\centerline{\sc II тур. \ 7 класс.}
\medskip
\hrule
\medskip

\q1. В тестировании приняло участие 200 школьников. Все их работы учитель разложил по нескольким пачкам, после чего собрал всех участников в большом зале и стал проверять пачки работ в некотором порядке. Проверка каждой работы занимает ровно одну минуту. Закончив проверять пачку работ, учитель мгновенно сообщает их результаты всем участникам и приступает к проверке следующей пачки. Каждый школьник подсчитал, какое время он провел в зале до оглашения своего результата. Докажите, что суммарное время ожидания не зависит от порядка, в котором учитель проверяет пачки.

\q2. В пятиугольнике $ABCDE$ углы при вершинах $B$, $C$, $D$ и~$E$ равны. Биссектриса угла $A$ пересекает сторону $CD$ в точке $F$. Докажите, что $AB+DF = AE+CF$. 

\q3. Существуют ли такие натуральные числа $a$ и $b$, что $a^2>b$ и 
$$
a+\text{НОД}(a^2-b, a-1) = b+ \text{НОД}(b^2+a, b+1)\ ?
$$

\q4. У аптекаря есть $n$ гирь, они покрашены в три цвета --- белый, синий и красный. 
Суммарный вес всех гирь каждого цвета равен 1 кг. %одному килограмму. 
Аптекарю нравится набор гирь, если его суммарный вес не меньше 1 кг. Докажите, что ему нравится хотя бы $\frac{3}{4} \cdot 2^n$ наборов.

\dotline

\centerline{Олимпиада 2024 года. II тур. 7 класс. Выводная аудитория.}
\smallskip

\q5. Дан треугольник $ABC$, в котором $\angle A = 2 \angle C$. На продолжении стороны $AB$ за точку $B$ отмечена такая точка $D$, что $AC = BD$. Докажите, что $AB + BC \geqslant CD$.

\q6. У Саши есть две бесконечные ленты. На первой ленте записаны числа 
$$
1, \quad 2, \quad 4, \quad 8, \quad 16, \quad 32, \quad 64, \quad 128, \quad \ldots
$$
(каждое следующее вдвое больше предыдущего). На второй ленте записаны числа: 
$$
1, \quad 2, \quad 4, \quad 8, \quad 61, \quad 23, \quad 46, \quad 821, \quad \ldots
$$
(<<перевернутые>> числа с первой ленты). На какую цифру может начинаться натуральное число $n$, если число $n^2$ записано на второй {ленте?\parfillskip=0pt\par}

\q7. В клетчатом квадрате $N \times N$ отмечены центры $3N$ клеток. Докажите, что среди попарных расстояний между отмеченными точками какие-то два отличаются в два раза.


\clearpage

\centerline{\sc САНКТ-ПЕТЕРБУРГСКАЯ}
\centerline{\sc ОЛИМПИАДА ШКОЛЬНИКОВ 2024 года ПО МАТЕМАТИКЕ}
\centerline{\sc II тур. \ 8 класс.}
\medskip
\hrule
\bigskip

\q1. Докажите, что хотя бы одно из следующих четырёх чисел 
\begin{center}
{\sc  икс, \ микс, \ пятиклассник, \ пляскитиматинататами}
\end{center}
составное (разным буквами соответствуют разные цифры, 
одинаковым~--- одинаковые, первые цифры чисел не равны 0). 

\q2. На каждой клетке доски $8\times 8$ стоит отличник или двоечник. Один школьник \emph{видит} другого, если стоит с ним на одной вертикали или одной горизонтали. У каждого из них спросили: \emph{Сколько отличников ты видишь?} Отличники ответили верно, а каждый двоечник ошибся ровно на~1, ничуть не беспокоясь о каком бы то ни было 
правдоподобии своего ответа. Мог ли каждый из ответов прозвучать ровно четыре раза?

\q3. У аптекаря есть $n$ гирь, они покрашены в три цвета --- белый, синий и красный. 
Суммарный вес гирь каждого цвета равен 1 кг. %одному килограмму.
Аптекарю \emph{нравится} набор гирь, если его суммарный вес не меньше 1 кг. %килограмма.
Докажите, что ему нравится хотя бы $\frac34\cdot 2^n$ наборов. %23-66

\q4. В выпуклом пятиугольнике $ABCDE$ диагональ $BE$ параллельна стороне 
$CD$. Кроме того, $\angle ABE=\angle ADB$ и $\angle ACE=\angle AEB$. 
Докажите, что точка $A$ равноудалена от прямых $BD$ и~$CE$.

\dotline

\centerline{Олимпиада 2024 года. II тур. 8 класс. Выводная аудитория.}
\smallskip

\q5. Город с населением в миллион человек
называется \emph{атомизированным}, если среди любых 1000 жителей
не более чем $998$ пар знакомых. Каково наибольшее
возможное общее количество пар знакомых в 
атомизированном городе? 

\q6. Внутри равнобедренного треугольника $ABC$ ($AB=AC$) отмечена такая точка $D$, что $\angle ABD = \frac12 \angle BCD$. Прямая $\ell$ проходит через точку~$D$ параллельно прямой $BC$ и пересекает треугольник $ABC$ по отрезку длины $a$. Докажите, что $CD>a$. 

\q7. Последовательность $x_1,x_2,\ldots$, состоящая из нечётных натуральных чисел, 
такова, что $x_n+n$ кратно $x_{n+1}$
при всех $n$. Обязательно ли среди членов последовательности есть единица?

\end{document}
